\section{Фактормножества и~факторпространства}

	\begin{definition}
		\label{def:pair}
		\tdef{Упорядоченная пара}~$(a, b)$ элементов%
		\footnote{Которые в~рамках теории множеств тоже считаются множествами.}%
		~$a$ и~$b$~"--- это множество
		\footnote{%
			Определение упорядоченной пары, как и~определение декартова
			произведения, может отличаться в~зависимости от того, с~какого взгляда
			на мир (теоретико"=множественного или какого"=нибудь ещё) мы стартуем.

			В~прошлом веке математики определили <<первичным>> понятием, к~которому
			сводились все остальные понятия, \emph{множество}. Можно было бы вместо
			него взять какое-нибудь другое.
		}%
		\[
			(a, b) := \set*{a, \set*{a, b}}
			.
		\]
	\end{definition}
	

	\begin{definition}
		\label{def:Decartes_product}
		\tdef{Декартовым произведением} множеств~$A$ и~$B$ называется множество
		упорядоченных пар их элементов:
		\[
			A \times B :=
			\set{(a, b)}{(a \in A) \wedge (b \in B)}
			.
		\]
	\end{definition}

	\begin{definition}
		\label{def:relation}
		Пусть $A$ и~$B$~"--- множества. \tdef{Отношением~$R$ на паре
		множеств~$A$ и~$B$} называется подмножество~$R \subset A \times B$ их
		декартова произведения.

		В~случае, если~$B = A$, то говорят об \tdef{отношении на
		множестве~$A$}.

		Говорят, что элементы~$a \in A$ и~$b \in B$ \tdef{находятся
		в~отношении~$R$} и~пишут $a \mathrel{R} b$, если $(a, b) \in R$:
		\[
			\bigl(
				a \mathrel{R} b
			\bigr)
			:=
			\bigl(
				(a, b) \in R
			\bigr)
			.
		\]
	\end{definition}

	Определение отношения~"--- это обобщение понятия графика функции.

	\begin{definition}
		Отношение~$\sim$ на множестве~$A$ называется \tdef{отношением
		эквивалентности}, если для всех~$a$, $b$ и~$c \in A$ выполнены условия:
		\begin{enumerate}[1)]
			\item \tdef{рефлексивность}: для любого $a \in A$ выполнено
				$a \sim a$;
			\item \tdef{симметричность}: для любых $a$, $b \in A$ выполнено
				\[
					(a \sim b) \implies (b \sim a);
				\]
			\item \tdef{транзитивность}: для любых $a$, $b$ и~$c \in A$ выполнено
				\[
					\bigl(
						(a \sim b)
						\wedge
						(b \sim c)
					\bigr)
					\implies
					\bigl(
						a \sim c
					\bigr)
					.
				\]
		\end{enumerate}
	\end{definition}

	Пусть $(A, \sim)$~"--- множество с~заданным на нём отношением
	эквивалентности.
	\begin{definition}
		\label{def:equivalence_class}
		\tdef{Класс эквивалентности~$[a]$} элемента~$a \in A$~"--- это множество
		всех элементов~$b \in A$, эквивалентных~$a$:
		\[
			[a] :=
			\set{b \in A}{b \sim a}
			.
		\]
	\end{definition}

	\begin{lemma}
		 Отношение эквивалентности разбивает множество~$A$ на непересекающиеся
		 подмножества, образующие классы эквивалентности.
	\end{lemma}
	\begin{proof}
		Пусть~$[a] \subset A$~"--- множество всех элементов,
		эквивалентных~$a \in A$, а~$[b] \subset A$~"--- множество всех
		элементов,
		эквивалентных~$b \in A$. Предположим, что $[a] \cap [b]$ не пусто.
		Значит, оно содержит в~себе хотя бы один элемент~$c \in [a] \cap [b]$.
		Теперь, так как~$c \in [a]$, то~$c \sim a$; а~так как $c \in [b]$, то
		$c \sim b$. Поэтому
		\[
			a \sim c \sim b~,
		\]
		откуда $a \sim b$. Но так как~$[a]$~"--- это множество всех элементов,
		эквивалентных~$a$, то возьмём из него произвольный%
		\footnote{%
			Возможно, совпадающий с~$a$.%
		}
		элемент~$a' \in [a]$; точно так же возьмём элемент%
		\footnote{%
			Возможно, совпадающий с~$b$.%
		}
		~$b' \in [b]$. Тогда%
		% для произвольных~$a' \in [a]$ и~$b' \in [b]$ имеем:
		\[
			a' \sim a \sim c \sim b \sim b'
			,
		\]
		откуда~$a' \sim b'$ для произвольных~"--- а~значит, и~всех~"---
		элементов~$a'$ и~$b'$. Так как множества~$[a]$ и~$[b]$
		состоят из одинаковых элементов, то они совпадают.%\ignorespaces
	\end{proof}

	\begin{definition}
		\label{def:factorset}
		Множество классов эквивалентности по отношению эквивалентности~$\sim$
		называется \tdef{фактормножеством}:
		\[
			A/{\sim} := \set{[a]}{a \in A}.
		\]
	\end{definition}

	\subsection{Пример фактормножества: конструирование множества целых чисел}

		\newcommand*{\tW}{\widetilde{W}}

Пусть~$A$ --- частично упорядоченное множество.

Определим функцию~$S$ её действием на множествах: $S(A) := A \cup \set*{A}$. Назовём множество~$S(A)$ множеством, \emph{следующим за~$A$}. Тогда

\begin{align*}
	0 &:= \emptyset,
	\\
	1 &:=
	S(0) =
	S(\emptyset) =
	\emptyset \cup \set*{\emptyset} =
	\set*{\emptyset},
	\\
	2 &:= S(1) =
	S \bigl( S(\emptyset) \bigr) =
	\set*{\emptyset} \cup \set*{\set*{\emptyset}} =
	\set*{\emptyset, \set*{\emptyset}},
	\\
	3 &:= S(2) =
	S \Bigl( S \bigl( S(\emptyset) \bigr) \Bigr) =
	\set*{\emptyset, \set*{\emptyset}}
	\cup
	\set*{\set*{\emptyset, \set*{\emptyset}}}
	=
	\set*{
		\emptyset, \set*{\emptyset}
		\set*{\emptyset, \set*{\emptyset}}
	}
	,
\end{align*}
и так далее.

Сложение чисел определяется тогда как объединение множеств:
\[
	1 + 2
	:=
	\set{1, 2} =
	\set*{\emptyset} \cup
	\set*{\emptyset, \set*{\emptyset}}
	=
\]

У нас есть конструкция натуральных чисел, начиная \emph{с нуля} (ноль мы считаем натуральным числом, следуя Н.\,Бурбаки). Теперь определим множество целых чисел~$\ZZ$.

Интуитивно хочется получить что-то, состоящее из нуля~$\set*{0}$, положительных натуральных чисел~$\NN \setminus \set*{0}$ и <<отрицательных натуральных чисел>> $-(\NN\setminus\set*{0}) = \set{-n}{n \in \NN \setminus \{0\}}$. Проблема в том, что операция <<минус>> не определена. Можно её определить аксиоматически, но мы пойдём другим путём.

Определим целые числа с помощью пар натуральных чисел. Геометрическая иллюстрация этого процесса показана на~\picref{pic:integers}: мы отождествляем точки вдоль прямых в~$\NN^2$ и <<нанизываем их>> на горизонтальный и вертикальный лучи --- экземпляры~$\NN$. Развернув <<угол>>, получившийся после отождествления точек, получим прямую, на которой живут целые числа.

Две точки $(a, b)$ и $(m, n)$ из~$\NN^2$ лежат на одной из прямых, изображённых на~\picref{pic:integers}, тогда и только тогда, когда $a - m = b - n$ (вектор с координатами $(a - m, b - n)$ коллинеарен направляющему вектору прямой).


Проблема только в том, что для натуральных чисел не определена операция вычитания.

Эта трудность не преодолевается, но обходится: заметим, что равенство $a - m = b - n$ равносильно равенству $a + n = b + m$. Сложение определено для натуральных чисел, так что его и примем за определение того, что $(a, b)$ и $(m, n)$ лежат на одной прямой. Поэтому
$\ZZ = \NN^2 / \sim$, где
\[
	\bigl(
		(a, b) \sim
		(m, n)
	\bigr)
	:=
	\bigl(
		a + n = b + m
	\bigr)
	.
\]
Обозначим класс эквивалентности числа~$(m, n) \in \NN^2$ символом $[(m, n)]$:
\[
	[(m, n)] :=
	\set{(a, b) \in \NN^2}{(a, b) \sim (m, n)}.
\]

При этом положительные числа в~$\ZZ$ эквивалентны точкам на горизонтальной прямой, а отрицательные --- точкам на вертикальной прямой.

% =============================================================================

Заметим, что при $a \in \NN$ и $k \in \NN$
\begin{align*}
	[(a, 0)] &= [(a + k, k)];
	\\
	[(0, a)] &= [(k, a + k)].
\end{align*}
Поэтому если в паре~$(a, b)$ имеем $a > b$, то класс $[(a, b)]$ представляет положительное число, а если $a < b$ --- то отрицательное.


Определим отображение $+ \colon \ZZ \times \ZZ \to \ZZ$. Если, интуитивно, при $a \in \NN$ класс $[(a, 0)] \in \ZZ$ определяет неотрицательное число $a \in \ZZ$, а класс $[(0, a)] \in \ZZ$ --- отрицательное число $-a \in \ZZ$, то определяем


Пока мы не ввели сложения и умножения в~$\NN^2 / \sim$. Подумаем, как можно его ввести, используя знания из школы.

\subsection[Вводим операции в $\NN^2/\sim$]{Вводим операции в $\bm{\NN^2/\sim}$}

\emph{Интуитивно} мы можем понять, что такое число $m - n$ при $m \geq n$. Поэтому при $m \geq n$ имеем $(m, n) \sim (m - n, 0)$ --- и пара $(m, n)$ представляет число $m - n \in \NN$. А при $m < n$ имеем $(m, n) \sim (0, n - m)$.

% =============================================================================

В этом разделе мы опишем аналог векторного пространства (или модуля над кольцом) для $\NN^2$. Но сначала надо с ним познакомиться.

\subsection{Определения множеств с операциями}

\begin{definition}
	Множество~$M$ с бинарной операцией $* \colon M \times M \to M$ называется \tdef{магмой}.
	\label{<+label+>}
\end{definition}
От операции~$*$ требуется только \tdef{замкнутость}: применение операции~$*$ к паре $(m, n) \in M \times M$ должно давать элемент множества~$M$.

\begin{definition}
	Если в магме~$(M, *)$ потребовать от операции~$*$ ещё и ассоциативности: $(a
	* b) * c = a * (b * c)$ для всех $a$, $b$, $c$ из~$M$, --- то получим
	определение \tdef{полугруппы}%
	\footnote{%
		Заметка на полях: как мы узнаем позже, $(\ZZ, +)$ --- это \emph{группа}.
		А~вот если считать, что в~$\NN$ нет нуля, то $(\NN, +)$ ---
		\emph{полу}группа, а её подлежащее множество~$\NN$--- это <<половина
		от~$\ZZ$>>. Возможно, поэтому дали такое название?
	}%
	.
\end{definition}

\begin{definition}
	Если в полугруппе $(M, *)$ есть элемент~$e \in M$, такой, что для любого $a \in M$ выполняются равенства $a * e = e * a = a$, то элемент~$e$ называется \tdef{нейтральным}, а $(M, *)$ --- \tdef{моноидом}.
	\label{<+label+>}
\end{definition}

Итак, множество~$\NN$ с операцией~$+ \colon \NN \times \NN \to \NN$ --- это моноид: сложение ассоциативно, а в качестве нейтрального (по сложению) элемента выступает~$0 \in \NN$.

Далее, множество~$\NN$ с операцией $\cdot \colon \NN \times \NN \to \NN$ --- это тоже моноид: умножение ассоциативно, а нейтральный (по умножению) элемент --- это $1 \in \NN$.

\subsubsection{Группа}
\begin{definition}
	Множество~$(G, *)$ c бинарной операцией $* \colon G \times G \to G$ называется \tdef{группой}, если:
	\begin{enumerate}[1)]
		\setcounter{enumi}{-1}
		\item $G$ замкнуто относительно~$*$;
		\item $*$ ассоциативна: $a * (b * c) = (a * b) * c$;
		\item в~$G$ существует нейтральный относительно~$*$ элемент~$e \in G$, такой, что для любого элемента $a \in G$
			\[
				a * e = e * a = a
				;
			\]
		\item у всякого $a \in G$ существует обратный относительно~$*$ элемент~$\hat{a} \in G$, то есть такой, что
			\[
				\hat{a} * a = e = a * \hat{a}
				.
			\]
	\end{enumerate}
	Если сложение удовлетворяет ещё аксиоме
	\begin{enumerate}
		\setcounter{enumi}{3}
		\item для всяких~$a$, $b \in G$ операция~$*$ коммутативна:
			\[
				a * b = b * a,
			\]
	\end{enumerate}
	то группа называется \tdef{абелевой}.
	\label{<+label+>}
\end{definition}

\subsubsection{Поле}

\begin{definition}
	\tdef{Поле} --- это тройка~$(F, +, \cdot)$, состоящая из множества~$F$ и двух
	операций~$+$ и~$\cdot$ на нём, удовлетворяющая условиям:
	\begin{enumerate}[1)]
		\item $(F, +)$ --- абелева группа с нейтральным элементом~$0$;
		\item $(F\setminus\set*{0}, \cdot)$ --- абелева группа с нейтральным элементом~$1$;
		\item $0 \cdot 1 = 0$;
		\item $a \cdot (b + c) = a\cdot b + a \cdot c$ для любых $a$, $b$ и $c \in F$;
	\end{enumerate}

	\label{<+label+>}
\end{definition}<++>

\subsubsection{Кольцо}

\subsubsection{Модуль над кольцом}

\subsection[Аналог модуля над кольцом для $\NN^2$]{Аналог модуля над кольцом для $\bm{\NN^2}$}

Опишем теперь

\subsubsection{Операции в~$\NN^2$}

Будем рассматривать $W$ как <<модуль>> над полукольцом~$\NN$ (в отличие от <<настоящего>> модуля над кольцом, у нас нет операции $v \mapsto -v$). Пусть $\bas{e} = (\bm{e}_1, \bm{e}_2)$ --- базис в~$W$. Будем считать, что $\bm{e}_1$ направлен вправо, а $\bm{e}_2$ --- вверх, как на~\picref{pic:integers}.

Будем отождествлять элемент $W \ni a\bm{e}_1 + b\bm{e}_2$ с набором координат~$(a, b) \in \NN^2$ в базисе~$\bas{e}$. Поэтому $W$ изоморфно $\NN^2$: изоморфизм устанавливается соответствием
\[
	\NN^2 \ni (a, b) \mapsto a\bm{e}_1 + b \bm{e}_2 \in W
	.
\]
Будем обозначать тот факт, что модуль $W$ изоморфен~$\NN^2$, символом $W \cong \NN^2$.

Сначала определим сложение в~$W \cong \NN^2$: пусть 
\[
	+ \colon \NN^2 \times \NN^2 \to \NN^2
	\colon
	\bigl(
		(a, b),
		(m, n)
	\bigr)
	\mapsto
	\bigl(
		a + m, b + n
	\bigr)
	.
\]
Мотивация для этого --- посмотреть на сложение в~$W$:
\[
	(a\bm{e}_1 + b\bm{e}_2) + (m\bm{e}_1 + n\bm{e}_2) =
	(a + m)\bm{e}_1 + (b + n)\bm{e}_2
	.
\]

Если мы хотим, чтобы на~\picref{pic:integers} вектор~$\bm{e}_1$ соответствовал будущему числу~$1 \in \ZZ$, а вектор --- будущему числу~$-1 \in \ZZ$, то достаточно положить их сумму эквивалентной~$0$: $\bm{e}_1 + \bm{e}_2 \sim 0$. Так как $0 = 0\bm{e}_1 + 0\bm{e}_2$, то в координатах это означает, что $(1, 0) + (0, 1) \sim (0, 0)$.

Теперь перейдём к фактормодулю $\tW := W/\sim$, изоморфному $\NN^2 / \sim =: \ZZ$.

Определим сложение в нём унаследованным из сложения в~$\NN^2$:
\[
	+\colon
	\tW
	\times
	\tW
	\to
	\tW
	\colon
	\bigl(
		[(a, b)],
		[(m, n)]
	\bigr)
	\mapsto
	[(a, b) + (m, n)]
\]
Иными словами, $[(a, b)] + [(m, n)] = [(a, b) + (m, n)]$. 

Для фактормодуля~$\tW$ поэтому имеем
\[
	[v] + [w] := [v+w]
\]
для любых векторов~$v$, $w \in W$.

Теперь определим билинейное умножение на~$\tW$, превратив его в <<алгебру>>. Если хотим уметь вычислять $[v] \cdot [w]$, то, разложив эти векторы по базису: $\bm{v} = v^1 \bm{e}_1 + v^2 \bm{e}_2$, $\bm{w} = w^1 \bm{e}_1 + w^2 \bm{e}_2$, --- получим, что
\[
	[v] \cdot [w] =
	[v^1 \bm{e}_1 + v^2 \bm{e}_2]
	\cdot
	[w^1 \bm{e}_1 + w^2 \bm{e}_2]
	=
	\bigl(
		v^1 [\bm{e}_1] +
		v^2 [\bm{e}_2]
	\bigr)
	\cdot
	\bigl(
		w^1 [\bm{e}_1] +
		w^2 [\bm{e}_2]
	\bigr)
	=
	\bigl(
		v^1 w^1 \,
		[\bm{e}_1] \cdot
		[\bm{e}_1]
		+
		v^1 w^2 \,
		[\bm{e}_1] \cdot
		[\bm{e}_2]
		+
		v^2 w^1 \,
		[\bm{e}_2] \cdot
		[\bm{e}_1]
		+
		v^2 w^2 \,
		[\bm{e}_2] \cdot
		[\bm{e}_2]
	\bigr)
	.
\]
Итак, умножение векторов в~$\tW$ определяется умножением классов $[\bm{e}_\nu]$
базисных векторов.  Теперь, в соответствии с тем, что $[\bm{e}_1]$ --- это
вектор в направлении <<положительных>> чисел, а $[\bm{e}_2]$ --- в направлении
<<отрицательных>>, определим умножение классов базисных векторов:
\[
	\begin{array}{r@{{} = {}}l@{\quad}l}
		[\bm{e}_1] \cdot [\bm{e}_1] & [\bm{e}_1]
		& \text{(<<$1 \cdot 1 = 1$>>);}
		\\
		[\bm{e}_1] \cdot [\bm{e}_2] & [\bm{e}_2]
		& \text{(<<$1 \cdot (-1) = -1$>>);}
		\\
		[\bm{e}_2] \cdot [\bm{e}_1] & [\bm{e}_1]
		& \text{(<<$(-1) \cdot 1 = -1$>>);}
		\\
		[\bm{e}_2] \cdot [\bm{e}_2] & [\bm{e}_1]
		& \text{(<<$(-1) \cdot (-1) = 1$>>).}
	\end{array}
\]

Сопоставляя $\NN^2 / \sim \ni [(a, b)] \longleftrightarrow a[\bm{e}_1] + b [\bm{e}_2] \in \tW$, имеем:
\begin{align*}
	[(a, b)] \cdot [(m, n)] &=
	[a \bm{e}_1 + b \bm{e}_2] \cdot [m \bm{e}_1 + n \bm{e}_2]
	\nlinea{=}
	am [\bm{e}_1] \cdot [\bm{e}_1] + 
	an [\bm{e}_1] \cdot [\bm{e}_2] + 
	bm [\bm{e}_2] \cdot [\bm{e}_1] + 
	bn [\bm{e}_2] \cdot [\bm{e}_2]
	\nlinea{=}
	am [\bm{e}_1] + 
	an [\bm{e}_2] + 
	bm [\bm{e}_2] + 
	bn [\bm{e}_1]
	\nlinea{=}
	(am + bn) [\bm{e}_1] +
	(an + bm) [\bm{e}_2]
	.
\end{align*}
Это соответствует тому, что $[(a, b)] \cdot [(m, n)] = [(am + bm, an + bm)]$.


	\subsection{Факторпространства}
	\begin{definition}
		Пусть~$W$~"--- векторное пространство над полем~$\kk$, а~$V$~"--- его
		подпространство. Определим на~$W$ отношение эквивалентности~$\sim$ так:
		\[
			(w_1 \sim w_2) :=
			(w_2 - w_1 \in V).
		\]

		Определим~$W/V$ как множество
		\[
			W/V := W/{\sim} =
			\set{[w]}{w \in W},
		\]
		где
		\[
			[w] := \set{h \in W}{h \sim w} = \set{h \in W}{w - h \in V}
			.
		\]

		Если разность векторов~$w - h$ лежит в~одном подпространстве
		(<<плоскости>>), то интуитивно это означает, что конец вектора~$h$
		будет лежать в~подпространстве (<<плоскости>>), образованной
		параллельным переносом подпространства~$V$ <<в~конец вектора~$w$>>.

		Другая интерпретация класса эквивалентности~$[w]$ вектора~$w$~"--- это
			сдвиг подпространства~$V$ на вектор~$w$:
		\[
			[w] = \set{w + v}{v \in V} = w + V~\]
		(упражнение).
		\label{<+label+>}
	\end{definition}


	Чтобы сделать из множества~$W/V$ векторное пространство, нужно ввести
	операции сложения классов~$\mhp \colon (W/V) \times (W/V) \to (W/V)$
	и~умножения их на число~$\mhc \colon \kk \times (W/V) \to (W/V)$. Это
	делается так.

	\begin{definition}[сложение и~умножение в~$\bm{W/V}$]
		По определению
%				\begin{gather*}
%					\xymatrix{
%						\llap{$
%							\mhp\colon{}%
%						$}
%						(W/V) \times (W/V)
%						\ar[r]{}
%						&
%						W/V~%						\\
%						([w_1], [w_2])
%						\ar @{|->} [r]{}
%						&
%						[w_1] \mhp [w_2] := [w_1 + w_2];
%					}
%				\end{gather*}
		\[
			\begin{array}{r@{}r@{}c@{}r@{}c@{}l}
				\mhp
				\colon
				{}
				&
				(W/V) \times (W/V)
				&
				{} \to {}
				&
				\multicolumn{1}{@{}l@{}}{W/V}
				&&
				\\
				&
				([w_1], [w_2])
				&
				{} \mapsto {} 
				&
				[w_1] \mhp [w_2]
				&
				{} := {}
				&
				[w_1 + w_2];
				\\
				\\
				\mhc
				\colon
				{}
				&
				\kk \times (W/V) 
				&
				{} \to {}
				&
				\multicolumn{1}{@{}l@{}}{W/V}
				&&
				\\
				&
				(a, [w_2])
				&
				{} \mapsto {}
				&
				a \mhc [w_2]
				&
				{} := {}
				&
				[a \cdot w_2].
			\end{array}
		\]
		Иными словами, сложить классы~$[w_1]$ и~$[w_2]$~"--- то же самое, что
		взять класс суммы их представителей~$[w_1 + w_2]$; то же верно для
		умножения.
	\end{definition}

	Проверим корректность определения~"--- иными словами, покажем, что~$\mhp$
	и~$\mhc$~"--- действительно \emph{отображения}, а~не просто отношения, то
	есть всюду определены и~функциональны. Применительно к~нашей ситуации это
	означает, что нужно проверить следующие утверждения:
	\begin{itemize}
		\item сложение определено для всех пар классов векторов;
		\item умножение определено для всех пар, состоящих из числа и~класса
			вектора;
		\item обе операции не зависят от выбора представителя класса.
	\end{itemize}%

	\begin{proof}
		Первые два утверждения имеют место, так как отображение
		факторизации~$\pi \colon W \to W/V$ сюръективно, то есть в~каждом
		классе эквивалентности~$[w] \in W/V$, рассмотренном как подмножество
		в~$W$, есть хотя бы один вектор, а~значит, в~правые части формул для
		сложения и~умножения <<есть что подставить>>.

		Последнее утверждение проверим так. Будем понимать под выражением~$w +
		V$ множество всех векторов вида~$w + v$, где $v \in V$. Тогда $[w] =
		w~+ V$,~и
		\begin{align*}
			[w_1] \mhp [w_2] &=
			(w_1 + V) + (w_2 + V) 
			\nlinea{=}
			w_1 + w_2 + V + V 
			\nlinea{=}
			w_1 + w_2 + V 
			\nlinea{=}
			[w_1 + w_2]
			.
		\end{align*}
		(Здесь мы использовали не сложение векторов, а~сложение \emph{множеств}
		векторов, которое даже не определили. Смысл выражения~$V + V = V$
		в~том, что мы складываем все возможные пары векторов~$v_1$ и~$v_2$
		из~$V$ и~получаем вектор~$v_1 + v_2$, который живёт в~$V$. Это так,
		потому что~$V$~"--- подпространство.)

		Аналогично проверяется независимость выбора представителя класса для
		операции умножения на число:
		\begin{align*}
			a \mhc [w] &=
			a \cdot (w + V)
			\nlinea{=}
			aw + aV
			\nlinea{=}
			aw + V~\nlinea{=}
			[aw]
			.
		\end{align*}
		(Аналогично,~$aV := \set{av}{v \in V}$, и~$av \in V$ для любого~$v \in
		V$, так как~$V$~"--- подпространство.)%
	\end{proof}


	\begin{definition}
		\tdef{Аффинное пространство}~"--- это пара~$(A, V)$, состоящая из
		множества точек~$A$ вместе с~векторным пространством~$V$, причём для
		любых точек~$a$ и~$b$ из~$A$ определён вектор~$\vec{ab} \in V$.
		\label{<+label+>}
	\end{definition}

%
%%\footnote{%
%%	Напомним, что \tdef{аффинное пространство} над полем~$\kk$~"--- это
%%	векторное пространство над тем же полем, снабжённое операцией 
%
%%	Говоря неформально, это множество точек, которые можно отождествить с~%
%%	радикс-векторами с~началом в~точке~$O$ и~у~которых определена разность:
%%	если $a$ и~$b$~"--- точки, то $b - a$~"--- это вектор с~началом в~точке~$a$
%%	и~концом в~$b$.
%%}
%
%\textstars
%
%Пусть $\AAA^3$~"--- аффинное пространство над~$\RR$. Рассмотрим в~$\AAA^3$
%кристаллическую решётку.
%
%Если мы возьмём по одной точке в~каждой ячейке, то получим
%факторпространство по ячейке. Элементы этого факторпространства~"--- точки,
%каждая из которых соответствует ячейке~"--- образуют \tdef{трансляционную
%решётку} кристалла.
%
%Выберем четыре периодически повторяющихся атома решётки и~построим на них
%репер в~$\RR^3$. Тогда можно перейти к~рассмотрению факторпространства по
%решётке, образованной этим репером,~"--- <<нарезать>> пространство~$\RR^3$
%на параллелепипеды, отвечающие элементарным ячейкам. Если~$a_0$, $a_1$,
%$a_2$, $a_3$~"--- векторы из произвольного начала координат к точкам,
%которые можно сделать узлами решётки, то рассмотрим
%\begin{align*}
%	A &:=
%	a_0 + 
%	\opspan_{\ZZ} \set*{a_0 a_1,\ a_0 a_2,\ a_0 a_3}
%	\nlinea{=}
%	a_0 + 
%	\opspan_{\ZZ} \set*{a_1 - a_0,\ a_2 - a_0,\ a_3 - a_0}
%	\nlinea{=}
%	\opspan_{\ZZ} \set*{a_0 + \tau_1,\ a_0 + \tau_2,\ a_0 + \tau_3},
%\end{align*}
%где $\tau_\nu := a_\nu - a_0$,~"--- линейную оболочку векторов, исходящих
%из точки~$a_0$, с~целыми коэффициентами. Отождествим векторы, которые можно
%получить сдвигами на векторы, целые кратные $\tau_\nu$, $\nu \in
%\set*{1, 2, 3}$:
%\[
%	(v \sim w) := (v - w \in A)
%\]
%"--* и~составим факторпространство
%\begin{gather*}
%	\RR^3/A~=
%	\set{[v]}{v \in \RR^3},
%	\\
%	[v] := \set{w \in \RR^3}{w \sim v} = \set{v + h}{h \in A} = v + A,
%\end{gather*}
%в~котором живут классы эквивалентности~$[v]$ векторов $v \in \RR^3$~"---
%иными словами, $[v]$~"--- это всевозможные сдвиги вектора~$v$ на векторы~$h
%\in A$%
%%, то есть $[v]$~"--- это сдвиг подпространства~$A$ на вектор~$v$,~"---
%%а~классы считаются эквивалентными, если разность их представителей лежит
%%в~$A$:
%%\[
%%	([v] = [w]) := ([v~"--- w] = [0]) :=
%%	(\exists h \in A \suchthat v~"--- v = h).
%%	%\\
%%	%(\exists h \in \left<a_0, a_1, a_2, a_3\right> \suchthat v~"--- w = h).
%%\]
%.
%
%
%Все узлы трансляционной решётки эквивалентны. Если через любые два узла
%провести прямую, то на этой прямой будет расположено бесконечно много
%других узлов, расположенных на равных расстояниях друг от друга. Эта прямая
%называется \tdef{узловой прямой}.
%
%%(Иначе можно сказать, что узловая прямая~"--- прямая в~$\RR^3/A$, проведённая вдоль одного из кусочков границы.)
%
%Сдвиг на расстояние между соседними узлами называется \tdef{трансляцией}
%вдоль данной узловой прямой.
%
%Проведём плоскость через две узловые прямые, проходящие через общий
%узел,~"--- или, иными словами, две прямые в~$\RR^3/A$. На этой плоскости
%расоположено бесконечно много других узловых прямых~"--- или, что то же
%самое, на этих прямых в~$\RR^3$ лежат только две точки из~$\RR^3/A$,
%у~которых счётное число прообразов в~$\RR^3$.
%
%Расстояние между двумя соседними плоскостями системы называется
%\tdef{межплоскостным расстоянием} и~обозначается~$d$.
%
%Введём в~трансляционной решётке систему координат, поместив её начало
%в~произвольный узел решётки и~направив её оси вдоль базиса, построенного на
%системе из трёх линейно независимых векторов, концы которых совпадают
%с~узлами решётки. Параллелепипед, натянутый на базисные векторы, называется
%\tdef{элементарной ячейкой}.
%
%Факторпространство~$\RR^3/A$~"--- это элементарная ячейка.
%
%\paragraph{Классификация элементарных ячеек:}
%\begin{itemize}
%	\item если внутри элементарной ячейки нет узлов, то ячейка называется
%		\tdef{примитивной};
%
%	\item если внутри элементарной ячейки нет узлов, но они есть в~центрах
%		граней, а~сама ячейка симметрична относительно плоскости, ортогональной
%		прямой, соединяющей узлы в~противоположных гранях, то такая ячейка
%		называется \tdef{базоцентрированной};
%
%	\item если внутри элементарной ячейки есть один узел в~центре, а~на
%		гранях нет, то такая ячейка называется \tdef{объёмноцентрированной};
%
%	\item если есть узел в~центре ячейки и~на центрах всех граней, то такая
%		ячейка называется \tdef{гранецентрированной}.
%\end{itemize}
%
%\paragraph{Индексы Миллера.}
%	Рассмотрим плоскость, не проходящую через начало координат, и~запишем её
%	уравнение в~отрезках:
%	\[
%		\frac{x}{A} +
%		\frac{y}{B} +
%		\frac{z}{C}
%		=
%		1,
%	\]
%	где $A$, $B$, $C$~"--- координаты точек на осях $Ox$, $Oy$ и~$Oz$
%	соответственно, в~которых плоскость пересекает эти оси.
%
%	Пусть $a$, $b$ и~$c$~"--- длины векторов трансляции вдоль осей $Ox$, $Oy$
%	и~$Oz$ соответственно. Длины отрезков~$A$, $B$ и~$C$, измеренные в~новых
%	единицах, будут выражаться рациональными числами~$p$, $q$, $r \in \QQ$%
%	\footnote{Почему?}:
%	\[
%		p = \frac{A}{a},
%		\qquad
%		q = \frac{B}{b},
%		\qquad
%		r = \frac{C}{c},
%	\]
%	так что можно найти такое целое число~$n$, что
%	\[
%		n = ph = qk = r \ell,
%	\]
%	где $h$, $k$, $\ell$~"--- целые
%	%, причём
%	%\[
%	%	h = \frac{n}{p},
%	%	\quad
%	%	k = \frac{n}{q},
%	%	\quad
%	%	\ell = \frac{n}{r}
%	%\]
%	и~$h \perp k \perp \ell$%
%	\footnote{%
%		Запись $a \perp b$ означает, что $a$ и~$b$
%		взаимно просты, то есть что $\gcd{a}{b} = 1$.%
%	}%
%	. 
%
%\begin{proof}
%	Достаточно взять~$h$, $k$ и~$\ell$ подходящими кратными числителей
%	чисел~$a$, $b$ и~$c$.\ignorespaces
%\end{proof}
%
%Набор чисел~$(h, k, \ell)$ называется \tdef{индексами Миллера} и
%обозначается~$[h, k, \ell]$.
%
%\subparagraph{Второй способ нахождения индексов Миллера.}
%В~прежних обозначениях
%\[
%	h := \frac{a}{A} \cdot q,
%	\qquad
%	k := \frac{b}{B} \cdot q,
%	\qquad
%	\ell := \frac{c}{C} \cdot q,
%\]
%где
%\[
%	q := \nok{\frac{A}{a},\ \frac{B}{b},\ \frac{C}{c}}
%	.
%\]
%
%
%Итак, уравнение плоскости приобретает вид
%\[
%	\frac{x}{h} + \frac{y}{k} + \frac{z}{\ell} = 1
%	.
%\]
%% Задаваемая этим уравнением плоскость, вообще говоря, не совпадает
%% с~исходной, однако получается сдвигом вдоль нормали исходной плоскости.
%% Эта плоскость проходит по осям~$Ox$, $Oy$ и~$Oz$ через узлы решётки.
%% Вдоль оси~$Ox$ нужно взять узел с~номером~$h$, вдоль $Oy$~"---
%% с~номером~$k$, вдоль $Oz$~"--- с~номером~$\ell$.
%
%\subparagraph{А вот ещё одно объяснение.}
%
%Рассмотрим уравнение плоскости в отрезках
%\begin{equation}
%\label{eq:Miller:ABC}
%\frac{x}{A} +
%\frac{y}{B} +
%\frac{z}{C}
%=
%1
%.
%\end{equation}
%Пусть~$A$, $B$ и~$C$ измеряются в единицах трансляции решётки, то есть
%выражаются рациональными числами%
%\footnote{%
%Если хотя бы одно из этих чисел не рационально, то прямая, стартовавшая из
%одного из узлов решётки, не пересечёт другого узла и поэтому не будет
%узловой прямой.%
%}%
%. Тогда можно найти число~$r \in \ZZ \setminus \set*{0}$, такое, что
%числа~$h := r/A$, $k := r/B$ и~$\ell := r/C$~"--- целые. Возьмём такое
%минимальное по модулю положительное число. Тогда
%уравнение~\eqref{eq:Miller:ABC} примет вид
%\begin{align*}
%\frac{1}{A}\,x +
%\frac{1}{B}\,y +
%\frac{1}{C}\,z
%&=
%1
%\\
%\frac{r}{A}\,x +
%\frac{r}{B}\,y +
%\frac{r}{C}\,z
%&=
%r
%\\
%hx + ky + \ell z &= r
%.
%\end{align*}
%Последнее уравнение~"--- каноническое уравнение плоскости, причём~$(h, k,
%\ell)$~"--- коэффициенты её вектора нормали.
%
%Итак, можно дать
%\begin{definition}
%\label{def:Miller_indeces}
%\tdef{Индексы Миллера}~"--- это наименьшие по модулю целочисленные
%коэффициенты вектора нормали плоскости, заданной в базисе, состоящем из
%векторов трансляции.
%\end{definition}
%
%\paragraph{Расчёт межплоскостного расстояния через индексы Миллера.}
%
%% Пусть $\bm{w} := (w_1, w_2, w_3)$~"--- базис, соответствующей ячейке, в~$\RR^3/A$, а~% $\bm{e} := (e_1, e_2, e_3)$~"--- стандартный ортонормированный базис в~$\RR^3$.
%% Пусть~$a_\nu := \abs{w_\nu}$~"--- длина $\nu$"=го базисного вектора. Тогда
%% \[
%% 	\bm{w} = \bm{e}\,C_{\bm{e}\bm{w}},
%% \]
%% где $C_{\bm{e}\bm{w}}$~"--- матрица перехода от базиса~$\bm{e}$ к~% базису~$\bm{w}$, в~$\nu$"=ом столбце которой стоят координаты $w_\nu^i$, $i \in \set*{1, 2, 3}$
%% вектора~$w_\nu$ в~базисе~$\bm{e}$.
%
%% Найдём межплоскостное расстояние. Очевидно, что в~направлении вектора~$w_\nu$ межплоскостное расстояние равно
%% \begin{equation}
%% 	t(w_\nu) := \frac{\abs{(w_1, w_2, w_3)}}{\norm{\vpr{w_{\mu_1}}{w_{\mu_2}}}},
%% 	\label{eq:interplane_dist}
%% \end{equation}
%% где~$(\mu_1, \mu_2) := (1, 2, 3) \setminus \set*{\nu}$~"--- упорядоченный по возрастанию индексов набор чисел без~$\nu$,
%% $(w_1, w_2, w_3)$~"--- смешанное произведение векторов
%% (объём~$\RR/A$, делённый на площадь основания).
%
%% Посмотрим теперь на индексы Миллера. Чтобы вычислить их, сначала переведём
%% уравнение плоскости~$\pi$, проходящей через концы базисных векторов~$\bm{w}$, в~% уравнение плоскости в~отрезках в~базисе~$\bm{e}$.
%
%% Вектор нормали~$n$ к~плоскости~$\pi$ имеет в~базисе~$\bm{w}$ вид
%% \[
%% 	n := \vpr{w_2~"--- w_1}{w_3~"--- w_1}.
%% \]
%% В~базисе~$\bm{e}$ тогда будем иметь
%% \[
%% 	n =
%% 	\begin{determ}
%% 		e_1 & e_2 & e_3
%% 		\\
%% 		w_2^1~"--- w_1^1 &
%% 		w_2^2~"--- w_1^2 &
%% 		w_2^3~"--- w_1^3
%% 		\\
%% 		w_3^1~"--- w_1^1 &
%% 		w_3^2~"--- w_1^2 &
%% 		w_3^3~"--- w_1^3
%% 	\end{determ}
%% 	.
%% \]
%
%% В~качестве векторов, относительно которых строим нормаль, можно было
%% взять~$w_{\mu_1}~"--- w_\nu$ и~$w_{\mu_2}~"--- w_\nu$.
%
%% Каноническое уравнение плоскости имеет вид
%% \[
%% 	Ax + By + Cz + D = 0.
%% \]
%% Если~$D \neq 0$, то можно преобразовать его в~уравнение плоскоси в~% отрезках:
%% \[
%% 	\frac{A}{-D}\,x +
%% 	\frac{B}{-D}\,y +
%% 	\frac{C}{-D}\,z~% 	=
%% 	1.
%% \]
%% При этом коэффициенты~$A$, $B$, $C$~"--- это координаты вектора нормали,
%% а~$-D$~"--- это скалярное произведение вектора нормали с~какой-нибудь
%% точкой, лежащей в~плоскости. Возьмём~$-D := \spr{n}{w_\nu}$. Тогда
%% уравнение плоскости в~отрезках примет вид
%% \[
%% 	\frac{n^1}{\spr{n}{w_\nu}}\,x +
%% 	\frac{n^2}{\spr{n}{w_\nu}}\,y +
%% 	\frac{n^3}{\spr{n}{w_\nu}}\,z~% 	=
%% 	1.
%% \]
%% Теперь
%% \[
%% 	h :=
%% 	\abs{w_1} \cdot
%% 	\frac{\spr{n}{w_\nu}}{n^1}
%% 	\cdot q,
%% 	\qquad
%% 	k :=
%% 	\abs{w_2} \cdot
%% 	\frac{\spr{n}{w_\nu}}{n^2}
%% 	\cdot q,
%% 	\qquad
%% 	\ell :=
%% 	\abs{w_3} \cdot
%% 	\frac{\spr{n}{w_\nu}}{n^3}
%% 	\cdot q,
%% \]
%% где
%% \[
%% 	q =
%% 	\nok{\frac{n^1}{\spr{n}{w_\nu}}, \frac{n^2}{\spr{n}{w_\nu}}, \frac{n^3}{\spr{n}{w_\nu}}}.
%% \]
%
%% В~частном случае, когда~$w_i = e_i$, имеем $w_i^j = \delta_i^j$, откуда $n~% = \frac{2}{\sqrt{3}} \cdot (1, 1, 1)^\top$, $\spr{n}{w_\nu} =
%% \spr{n}{e_\nu} = n^\nu = 2/\sqrt{3}$, $A = B = C = a = b = c = 1$, $q = 1$,
%% и~% \[
%% 	t(w_\nu) =
%% 	\frac{\abs{(w_1, w_2, w_3)}}{\norm{\vpr{w_{\mu_1}}{w_{\mu_2}}}} =
%% 	\frac{1}{1 \cdot \sgn{\nu, \mu_1, \mu_2}}.
%% \]
%
%В~частном случае для кубической решётки~"--- $a = b = c = 1$ и~$\spr{w_i,
%w_j} = \delta_{ij}$~"--- имеем:
%\begin{equation}
%	d =
%	\frac{a}{\sqrt{h^2 + k^2 + \ell^2}}
%	.
%	\label{eq:dCube}
%\end{equation}

% vim:ts=2:sw=2

